% Nouveau EGA V (cote 25): archive cote 25, site page 87 (PDF page 88).
% typescript. Draft compared: none.
% AI-transcribed by Claude Opus 5.5 (claude-opus-5-5, adaptive thinking, effort high), Message Batches API, batch ega25b (msgbatch_01A7n5YuXubUPNK1zsepi8RG), 2026-11-13; checked against both supplied scan parts; see CHECK_DECISIONS.tsv.
\scanpage{25/87}
\mspage{81}

suivant qu'on suppose  dim X \hand{$\geqslant$} m , ou  dim X \hand{$\leqslant$} m , voire dim X \hand{$<$} m.
Par la suite, nous regardons la $C_{\hand{\unsure{\eta|y}}}\hand{\subset} P_{k(\hand{\unsure{\eta|y}})}$ , correspondant au point
générique \hand{\unsure{\eta|y}} de  Grass$^{m+1}$ , et nous dispensons  de faire l'interprétation
des résultats obtenus en termes de  ``presque tous les points'' ...

Tout d'abord, nous avons déjà noté \var{\xout{\unsure{que|(cf}}}{}{typewriter strike-out}{ega25_25_087-v0001}dans 5.3. (et rattrapage du
cas général dans nº 12)  que $C_{\hand{\unsure{\eta|y}}}$  coupe ``régulièrement'' $X_{\hand{\unsure{\eta|y}}}$ , de façon
plus pré\unsure{c|s}ise et plus générale, pour tout F quasi-cohérent sur X,
la section $\hand{\phi}_{\hand{\unsure{\eta|y}}}^{\hand{(m+1)}}$ du \var{\inserted[typed]{Module loc libre}}{Module loc libre}{insertion}{ega25_25_087-v0002}\var{\xout{fibré vectoriel}}{}{typewriter strike-out}{ega25_25_087-v0003} de rang  m+1  sur  $X_{k(\hand{\unsure{\eta|y}})}$  dont le
schéma des zéros est \var{\over{C}{$C$}}{$C$}{writing over}{ega25_25_087-v0004}$_{\hand{\unsure{\eta|y}}}$ \var{\over{;}{\unsure{;|,}}}{\unsure{;|,}}{writing over}{ega25_25_087-v0005} est F-régulière. En vertu de 10.2. , cela
implique par exemple que le morphisme  $\overset{\hand{\sim}}{X}(\hand{C_{\unsure{\eta|y}}})\hand{\to} X_{k(\hand{\unsure{\eta|y}})}$ identifie $\overset{\hand{\sim}}{X}(\hand{C_{\unsure{\eta|y}}})$ au préschéma
déduit de $X_{k(\hand{\unsure{\eta|y}})}$ par éclatement de $f^{-1}(C_{\hand{\unsure{\eta|y}}})= Xx_P C_{\hand{\unsure{\eta|y}}}$ ;  \var{\xout{\unsure{D’autre part,}}}{}{typewriter strike-out}{ega25_25_087-v0006} dans le cas
où  dim f(X) \hand{$\leqslant$} m ,  on  aura même  $f^{-1}(C_{\hand{\unsure{\eta|y}}}) = \hand{\emptyset}$ , et  par suite
$\overset{\hand{\sim}}{X}(C_{\hand{\unsure{\eta|y}}}) \hand{\xrightarrow{\sim}} X_{k(\hand{\unsure{\eta|y}})}$  est un \uline{isomorphisme} , (et en fait, la \var{\over{p}{\unsure{p}}}{\unsure{p}}{writing over}{ega25_25_087-v0007}rojection conique
restreinte  est alors définie sur X tout entier à priori). Il se pose
ensuite la question de la dimension des fibres de $p_{C_{\hand{\unsure{\eta|y}}}}$ : $\overset{\hand{\sim}}{X}(C_{\hand{\unsure{\eta|y}}}) \hand{\unsure{\longrightarrow|\relbar}} \hand{Q}(C_{\hand{\unsure{\eta|y}}})$,
et de la platitude de ce morphisme. On trouve:
\uline{Proposition} 14.1.  \var{\xout{a) Si dim X}}{}{typewriter strike-out}{ega25_25_087-v0008} Supposons X irréductible,\var{\over{\unsure{e}}{\unsure{ou}}}{ \unsure{ou} }{writing over}{ega25_25_087-v0009}plus
généralement, que  pour toute composante irréductible $X_i$ de X, la fibre
de $X_i$ en le point $f(x_i)$  ($x_i$ = point générique de $X_i$) ait une dimension d
(indépendante de i), ce qui a lieu par exemple \var{\inserted{(avec d=0)}}{(avec d=0)}{insertion}{ega25_25_087-v0010} si f:X$\hand{\to}$P est quasi-fini\hand{,} \var{\margin{\var{\deleted{avec d=0}}{}{deletion}{ega25_25_087-v0012}}}{\placedmargin{\var{\deleted{avec d=0}}{}{deletion}{ega25_25_087-v0012}}}{marginal note placed}{ega25_25_087-v0011}
Alors :     a)  Si  dim f(X) \hand{$>$} m, alors la dimension  des fibres de
$p_C$  : $\overset{\hand{\sim}}{X}(C_{\hand{\unsure{\eta|y}}}) \hand{\longrightarrow} \var{\over{Q}{Q}}{Q}{writing over}{ega25_25_087-v0013}(C_{\hand{\unsure{\eta|y}}})$  sont toutes égales à  dim X - m .   b)  Si dim X
$\hand{\leqslant}$ m, \var{\inserted{et si les fibres non vides des $X_i$ sur P sont \unsure{ttes|tts} de dim d, alors;}}{et si les fibres non vides des $X_i$ sur P sont \unsure{ttes|tts} de dim d, alors;}{insertion}{ega25_25_087-v0014} les fibres \var{\inserted{non vides}}{non vides}{insertion}{ega25_25_087-v0015} de $p_C$ \var{\inserted{$\colon X_{k(\unsure{\eta|y})}\to Q(C_{\unsure{\eta|y}})$}}{$\colon X_{k(\unsure{\eta|y})}\to Q(C_{\unsure{\eta|y}})$}{insertion}{ega25_25_087-v0016} sont toutes de dimension \var{\over{0}{d}}{d}{writing over}{ega25_25_087-v0017}, \var{\deleted{i.e. $p_C$ est quasi-%
fini}}{}{deletion}{ega25_25_087-v0018} (donc $p_C$ est fini \var{\inserted{\unsure{resp. quasi-fini}}}{\unsure{resp. quasi-fini}}{insertion}{ega25_25_087-v0019} \var{\deleted{si X est propre sur \unsure{k$(\eta)$|k$(y)$} et}}{}{deletion}{ega25_25_087-v0020} \var{\inserted{$f\colon X\to P$ est fini \unsure{resp. quasi-fini}}}{$f\colon X\to P$ est fini \unsure{resp. quasi-fini}}{insertion}{ega25_25_087-v0021}\hand{]}.
% NOTE: the circled right-margin resp. quasi-fini is already encoded in the preceding insertion.

On \var{\deleted{se ramène aussitôt, en décomposant X en ses composantes irréd.,}}{}{deletion}{ega25_25_087-v0022}
\var{\deleted{au cas où X est irréductible}}{}{deletion}{ega25_25_087-v0023}. Dans le cas a), \var{\xout{si il y avait une fibre}}{}{typewriter strike-out}{ega25_25_087-v0024}
\var{\xout{correspondant à un}}{}{typewriter strike-out}{ega25_25_087-v0025} on sait déjà (je l'espère !) que pour \uline{tout}  point \hand{\unsure{\zeta|y}} de

\var{\margin{87}}{\placedmargin{87}}{marginal note placed}{ega25_25_087-v0026}
